\documentclass[aspectratio=169]{beamer}
\usetheme{Madrid}
\setbeamerfont{subtitle}{shape=\upshape}
\setbeamerfont{institute}{shape=\upshape}
\setbeamerfont{author}{shape=\upshape}
\setbeamerfont{date}{shape=\upshape}
\usepackage[T1]{fontenc}
\usepackage[utf8]{inputenc}
\usepackage[english]{babel}
\usepackage{booktabs}
\usepackage{amsmath}

\title{ML-SUPERB \& JEPA}
\subtitle{Trois slides pour le rapport principal (\texttt{report.tex})}
\author{Janis Vadim Bruny}
\date{3 avril 2026}

\begin{document}

\begin{frame}
\titlepage
\end{frame}

\begin{frame}{Problème et positionnement}
\begin{itemize}
\item Protocole \textbf{ML-SUPERB} : encodeur SSL \textbf{gelé}, petite tête \textbf{CTC}, budgets \textbf{10 min / 1 h} par langue.
\item Question centrale : comparer des représentations \textbf{JEPA / WavJEPA} à \textbf{HuBERT} sous la \textbf{même} recette downstream.
\item Résultats ASR : on rapporte \textbf{CER} et \textbf{WER} après décodage et scoring MLSUPERB — \textbf{tableau agrégé} sur la slide suivante, détail par langue dans \texttt{report.tex}.
\item Multilingue : mélange \texttt{eng} / \texttt{deu} / \texttt{fra} ; conditions comparées : \textbf{ASR seul} (10 min et 1 h), \textbf{ASR+LID} (10 min), \textbf{LoRA} sur ASR seul (10 min).
\end{itemize}
\end{frame}

\begin{frame}{Représentations et benchmarks (mono, multilingue, ABX)}
\small ABX mesure la séparabilité phonétique dans le latent ; CER/WER mesurent l’alignement texte.

\vspace{0.3em}
\begin{columns}[T]
\begin{column}{0.52\textwidth}
\textbf{Multilingue \& adaptation} \\
\small
\begin{tabular}{@{}lcc@{}}
\toprule
Réglage & CER & WER \\
\midrule
ASR-only 10 min & 24.96 & 23.48 \\
ASR-only 1 h & 20.76 & 18.30 \\
ASR+LID 10 min & 26.33 & 25.49 \\
LoRA ASR-only 10 min & 24.95 & 23.66 \\
\bottomrule
\end{tabular}

\vspace{0.2em}
\small
\begin{itemize}\setlength\itemsep{0.1em}
\item LoRA : pas de gain
\item 10 min $\rightarrow$ 1 h : grand écart
\end{itemize}
\end{column}
\begin{column}{0.46\textwidth}
\textbf{ABX} \\
\vspace{0.2em}
\begin{tabular}{@{}lc@{}}
\toprule
Modèle & err. ABX \\
\midrule
HuBERT & 0.5216 \\
WavJEPA HF & 0.5895 \\
$\Delta$ (W $-$ H) & +0.0679 \\
\bottomrule
\end{tabular}

\vspace{0.2em}
\small
\begin{itemize}\setlength\itemsep{0.1em}
\item HuBERT mieux sur ABX.
\item $\Delta$ ABX positif : HuBERT devant.
\end{itemize}
\end{column}
\end{columns}
\end{frame}

\begin{frame}{Dérive des poids}
\begin{columns}[T]
\begin{column}{0.55\textwidth}
\textbf{Temps de training} \\
\vspace{0.2em}
\begin{tabular}{@{}lc@{}}
\toprule
Piste & L4 (h) \\
\midrule
Multi.\ ASR 10 min & 5.93 \\
Multi.\ ASR 1 h & 11.22 \\
LID-only 10 min & 4.93 \\
LID-only 1 h & 11.17 \\
WavJEPA ASR \texttt{eng1} HF & 1.99 \\
WavJEPA ASR \texttt{eng1} local & 1.62 \\
\bottomrule
\end{tabular}
\end{column}
\begin{column}{0.43\textwidth}
\textbf{Dérive des poids} \\
\small
\begin{itemize}\setlength\itemsep{0.15em}
\item Cosinus moyen (HF / local) $\approx -0{,}008$ ; ratio normes local/HF $\approx 1{,}59$.
\item Lignes \texttt{ctc\_lo} : cosinus moyen $\approx 0{,}032$ entre runs HF et local.
\end{itemize}
\end{column}
\end{columns}

\vspace{0.35em}
\small \textbf{Conclusion :} la perte SSL peut baisser sans aligner le CER/WER ; espace latent et géométrie CTC peuvent \textbf{tourner} tout en gardant des scores ASR proches en bas budget.
\end{frame}

\end{document}
